%------------------------------------------------------------------------------%
\section{Apresentação e Motivação}

Apresentamos nesse capítulo o conceito de sequências de números reais.
Estamos interessados em explorar suas propriedades de convergência.
Esses resultados serão fundamentais quando iniciarmos a discussão sobre séries
no próximo capítulo.

Note que os mesmos conceitos podem ser estendidos para sequências
de outros objetos matemáticos, como vetores, matrizes ou funções.
Porém, as propriedades de convergência serão basicamente as mesmas.

Sequências e séries estão entre os conceitos matemáticos mais aplicados,
por exemplo, qualquer algoritmo computacional que busque a solução para
um problema fazendo correções iterativas em uma aproximação inicial vai produzir
uma sequência de aproximações que podem convergir ou não para a solução desejada.
Como exemplo, apresentamos o método de Newton.

%------------------------------------------------------------------------------%
\subsection*{Método das Tangentes de Newton}

Muitos métodos numéricos ou computacionais se baseiam na busca de uma solução
por aproximações sucessivas, isso é, eles partem de um chute inicial e
a cada iteração constroem uma nova aproximação para a solução.
Um dos métodos mais importantes com essas características é o
\emph{Método de Newton}\index{Método!Newton}\index{Newton!Método}
para encontrar zeros de funções.

Queremos encontrar um valor real $x^\ast$ que satisfaça a equação
\[
	f(x^\ast) = 0 \;,
\]
para uma função $f\colon A\subset\R\to\R$ derivável em $A$ e
cuja derivada, $f'$, é diferente de zero para todo $x\in A$.

O Método de Newton busca uma solução para $f(x)=0$ partindo de um chute inicial
$x_0\in A$. Começamos construindo a reta tangente ao gráfico de $f$ pelo ponto
$(x_0, f(x_0))$ e encontrando o ponto onde a reta toca o eixo-$x$, como ilustrado
na Figura~\ref{fig:newton_1}. O ponto $x_1$ encontrado dessa forma é a nova
aproximação para a solução.

%------------------------------------------------------------------------------%
\begin{figure}
  \centering
  \begin{subfigure}[b]{0.5\linewidth}
    \centering
    \input{fig_newton_step_1}
    \caption{Reta tangente em $(x_0, f(x_0))$\\e intercepto $x_1$}
    \label{fig:newton_1}
  \end{subfigure}%
  \begin{subfigure}[b]{0.5\linewidth}
    \centering
    % Title: gl2ps_renderer figure
% Creator: GL2PS 1.4.2, (C) 1999-2020 C. Geuzaine
% For: Octave
% CreationDate: Fri Sep 10 17:18:05 2021
\setlength{\unitlength}{1pt}
\begin{picture}(0,0)
\includegraphics[width=8cm,height=4.6cm]{fig_newton_step_2}
\end{picture}%
\begin{picture}(226,130)(0,0)
\fontsize{10}{0}\selectfont\put(53.553,61.3841){\makebox(0,0)[b]{\textcolor[rgb]{0,0,0}{{$x_0$}}}}
\fontsize{10}{0}\selectfont\put(169.103,51.1271){\makebox(0,0)[tl]{\textcolor[rgb]{0,0,0}{{$x_1$}}}}
\fontsize{10}{0}\selectfont\put(131.639,66.5127){\makebox(0,0)[br]{\textcolor[rgb]{0,0,0}{{$x_2$}}}}
\end{picture}

    \caption{Reta tangente em $(x_1, f(x_1))$\\e intercepto $x_2$}
    \label{fig:newton_2}
  \end{subfigure}
  \caption{Ilustração dos primeiros passos do Método de Newton}
\end{figure}
%------------------------------------------------------------------------------%

De posse da nova aproximação repetimos o processo para obter $x_2$, como ilustrado na
Figura~\ref{fig:newton_2}. Aplicando sucessivamente esse método produzimos uma
lista infinita de aproximações para a solução do problema.

Podemos construir uma fórmula para calcular cada aproximação $x_n$ a partir do valor da
aproximação anterior $x_{n-1}$. O primeiro passo é calcular a expressão para a reta
tangente ao gráfico de $f$ no ponto $(x_n, f(x_n))$
\[
	y-f(x_n) = f'(x_n)(x-x_n)    \;.
\]
A próxima aproximação será o valor de $x$ onde essa reta cruza o eixo-$x$, isso é,
quanto $y=0$, assim temos
\[
	-f(x_n) = f'(x_n)(x_{n+1}-x_n)      \;,
\]
rearranjando, obtemos a fórmula para o Método de Newton
\[
	x_{n+1} = x_n - \frac{\,f(x_n)\,}{\,f'(x_n)\,}    \;.
\]
Note que a exigência de que a derivada de $f$ não seja nula em $A$ é para
garantir que a divisão sempre seja possível.

Esse método vai produzir várias listas infinitas de números
\[
\begin{array}{llllll}
	x_0,     & x_1,     & x_2,     & \dots, & x_n,     & \dots \\
	f(x_0),  & f(x_1),  & f(x_2),  & \dots, & f(x_n),  & \dots \\
	f'(x_0), & f'(x_1), & f'(x_2), & \dots, & f'(x_n), & \dots
\end{array}
\]
a cada uma dessas listas infinitas damos o nome de \emph{Sequência Numérica}.

\OnlineBox{https://colab.research.google.com/drive/1li98CA9JUOYCO2_Gca8-Yspuex61ndH9?usp=share_link}
{
  Código Python para ilustrar o Método de Newton.
}

A pergunta que precisamos responder agora é se essas sequências
convergem para a solução que desejamos. O primeiro passo é definir
o que queremos dizer com convergência, depois verificamos se as sequências
convergem para alguma coisa e por fim verificamos que os valores
$f(x_n)$ convergem para zero e consequentemente os valores $x_n$
convergem para a solução que buscamos. Nas próximas seções
deste capítulo apresentamos os resultados matemáticos necessários
para respondermos a essas questões.

%------------------------------------------------------------------------------%
